\section*{Technical Validation}
%%% https://www.nature.com/sdata/submission-guidelines
% - describe the experiments, analyses or checks needed to support the technical quality of the dataset, with any supporting figures and tables, as needed.

The automated extraction pipeline includes two internal verification stages, evidence-span alignment and LLM-based factual verification, that discard unsupported extractions before inclusion in the final corpus (see Methods). To assess the quality of the extracted fields independently of these automated checks, we conduct a manual expert validation on a stratified sample of the corpus.

\subsection*{Sampling design}
The validation population comprises the 72,678 reviews with parsed full texts, i.e.\ all documents on which the extraction pipeline was executed. The unit of validation is the document: for each sampled review, all extracted methodological fields are assessed against the source PDF. Results are reported per field type.

The sample size is determined using the standard formula for estimating a proportion~\cite{cochran:1977}. At a 95\% confidence level ($z = 1.96$) and a margin of error of $\pm 5$ percentage points ($e = 0.05$), and using the conservative assumption $p = 0.5$, which maximizes the required sample size independently of the true extraction quality, the base sample size is
\[
n_0 = \frac{z^2 \, p(1-p)}{e^2} = \frac{1.96^2 \times 0.5 \times 0.5}{0.05^2} \approx 385.
\]
Applying the finite population correction for $N = 72{,}678$,
\[
n = \frac{n_0}{1 + \frac{n_0 - 1}{N}} \approx 383,
\]
which we round up to 385 documents. Because the correction is negligible at this population size, the sample size is effectively independent of $N$.

Documents are drawn by proportionate stratified random sampling, with strata defined by field and by document length. Stratification ensures that non-medical disciplines and longer documents, where parsing and extraction are expected to be more difficult, are represented according to their share of the population.

As not every review reports every methodological field, per-field sample sizes vary with field coverage: frequently reported fields (e.g.\ objective, date range) occur in most sampled documents, while sparsely reported fields (e.g.\ databases queried, research questions) occur in fewer. Estimates for sparsely covered fields are correspondingly less precise. We therefore report the exact number of judgments and the associated confidence interval for each field.

\subsection*{Annotation protocol}
Three annotators with expertise in systematic review methodology assess the sample against written annotation guidelines, fixed before annotation, that define correct and incorrect extractions and specify the treatment of boundary cases such as partial extractions and formatting variants. A shared subset of 80 documents is annotated independently by all three annotators to quantify inter-annotator agreement; the remaining 305 documents are divided evenly among them.

For each sampled document, two complementary assessments are performed. First, every non-null extracted field is compared against its verbatim evidence span and the source PDF, yielding a per-field correctness judgment (precision). Second, every field for which the pipeline returned no value is checked against the source PDF to determine whether the corresponding information is in fact reported, yielding an estimate of coverage (recall). The latter quantifies the pipeline's stated trade-off of completeness for precision.

Inter-annotator agreement on the shared subset is measured with Fleiss'~$\kappa$~\cite{fleiss:1971}; disagreements are resolved by discussion, and the adjudicated labels are used for the final estimates.

\subsection*{Reported estimates}
For each field type, precision and recall are reported together with 95\% Wilson score intervals~\cite{wilson:1927}, which remain accurate for proportions close to the boundaries~\cite{brown:2001}, the regime expected for a precision-oriented pipeline. Systematic error patterns observed during annotation are summarized in a qualitative error analysis.

% TODO 
% Sampling design (report the concrete sample composition)
% Annotation protocol
% Inter-annotator agreement
% Results Table
% Error analysis
